\documentclass[12pt,a4paper]{report}

% Pacotes básicos
\usepackage{mathptmx}
\usepackage{indentfirst}
\usepackage[utf8]{inputenc}    % Permite acentuação
\usepackage[T1]{fontenc}       % Codificação correta das fontes
\usepackage[brazil]{babel}     % Idioma português
\usepackage{geometry}          % Configuração da página
\geometry{a4paper, top=2cm, bottom=2cm, left=2.5cm, right=2.5cm}
\usepackage{setspace}          % Espaçamento entre linhas
\usepackage{graphicx}          % Inserir imagens
\usepackage{hyperref}          % Links clicáveis
\usepackage{titlesec}          % Personalização dos títulos
\usepackage{array}             % Tabelas mais flexíveis
\usepackage{longtable}         % Tabelas longas
\usepackage{tikz}
\usetikzlibrary{positioning,shapes.geometric}

% Configurações do documento
\onehalfspacing               % Espaçamento 1.5
\hypersetup{
    colorlinks=true,
    linkcolor=blue,
    urlcolor=blue,
    pdftitle={Universidade Estadual do Norte Fluminense Darcy Ribeiro},
    pdfauthor={Davi Rodrigues Soares Machado//Paulo Fernando Pereira Junior},
}

\begin{document}
\begin{titlepage}
    \begin{center}
        
        
\includegraphics[scale=0.2]{imagens/uenf1.png}\\[1cm]
       
        {\large CENTRO DE CIÊNCIAS E TECNOLOGIAS}\\
        {\large LCC - Laboratório de Ciência da Computação}\\[4cm]

        {\large \textbf{Davi Rodrigues Soares Machado}}\\
        {\large \textbf{Paulo Fernando Pereira Junior}}\\[4cm]

        {\large \textbf{Projeto da disciplina Paradigmas Orientados à Objetos para Desenvolvimento de Software: \\[2cm] Sistema de Aluguel de Carros}}\\[4cm]

        Campos dos Goytacazes - RJ\\
        \today
    \end{center}
\end{titlepage}


% Sumário
\tableofcontents
\newpage


% Introdução -----------------------------------------------------------------------
\chapter{Introdução}
O projeto consiste no desenvolvimento de um sistema para gerenciamento e
aluguel de veículos. A plataforma permitirá a interação entre clientes e agências de aluguel,
possibilitando a consulta de veículos, realização de reservas, pagamentos, acompanhamento do
aluguel e gerenciamento da frota.

Além das funcionalidades tradicionais, o projeto busca incorporar recursos de automação, como a identificação da quilometragem do veículo a partir de
fotografias, controle do estado do veículo, acompanhamento de custos operacionais e emissão
de alertas relacionados à manutenção e desvalorização.

O andamento do projeto estará disponível no \href{https://github.com/DaviRodrish/POODEV2026}{Repositório Github}.

% Objetivos -------------------------------------------------------------------------
\chapter{Objetivos}

\subsection*{Objetivo Geral}
Desenvolver um sistema para gerenciamento de uma agência de aluguel de veículos, permitindo controlar clientes, veículos da agência, reservas e pagamentos, manutenções e demais processos necessários para uma locação de veículo.
    

	\paragraph{Objetivos específicos}
    \begin{itemize}

    \item Permitir o cadastro e a autenticação de clientes,
    motoristas e agências;

    \item Permitir a consulta e reserva de veículos disponíveis;

    \item Diferenciar reservas de veículos novos e usados;

    \item Permitir o registro das condições do veículo por meio
    de fotografias;

    \item Registrar e controlar a quilometragem dos veículos;

    \item Detectar automaticamente a quilometragem apresentada
    em fotografias do painel do veículo;

    \item Permitir o registro do nível de combustível na retirada
    e na devolução;

    \item Disponibilizar pagamentos por PIX e cartão de crédito;

    \item Permitir a realização de entregas e devoluções fora da
    agência mediante cobrança adicional;

    \item Manter o histórico de aluguéis dos usuários e veículos;

    \item Disponibilizar informações relacionadas à manutenção
    dos veículos;

    \item Calcular o custo operacional dos veículos;

    \item Estimar a poluição gerada pela utilização dos veículos;

    \item Permitir o acionamento de um botão de emergência do seguro, enviando informações como localização e possíveis defeitos;

    \item Gerar alertas relacionados à desvalorização dos veículos;

    \item Aplicar regras de bloqueio de usuários de acordo com
    as condições definidas pela agência.

\end{itemize}

    
\chapter{Usuários do Sistema}

O sistema contará com diferentes perfis de usuários, cada um
com diferentes permissões e responsabilidades.

\section{Cliente / Locador}

O cliente poderá utilizar o sistema para consultar veículos,
realizar reservas, efetuar pagamentos e acompanhar seus aluguéis.

Entre suas principais funcionalidades estão:

\begin{itemize}

    \item Criar uma conta;
    \item Realizar login;
    \item Cadastrar informações pessoais;
    \item Enviar fotografia da CNH;
    \item Consultar veículos disponíveis;
    \item Realizar reservas;
    \item Efetuar pagamentos;
    \item Cadastrar motorista secundário;
    \item Consultar o histórico de aluguéis;
    \item Acionar o seguro em situações de emergência;
    \item Informar problemas encontrados no veículo.

\end{itemize}

\section{Agência}

A agência será responsável pelo gerenciamento da frota e dos
aluguéis.

Suas principais funcionalidades serão:

\begin{itemize}

    \item Cadastrar e gerenciar veículos;
    \item Consultar reservas;
    \item Gerenciar retiradas e devoluções;
    \item Registrar fotografias dos veículos;
    \item Registrar quilometragem e combustível;
    \item Acompanhar manutenções;
    \item Registrar problemas dos veículos;
    \item Consultar custos operacionais;
    \item Acompanhar a desvalorização dos veículos;
    \item Gerenciar usuários bloqueados;
    \item Acompanhar situações de emergência.

\end{itemize}


\chapter{Levantamento de Requisitos}

Durante a reunião inicial com o stakeholder, foram levantados
requisitos relacionados às funcionalidades de aluguel,
gerenciamento de veículos, pagamentos, manutenção e segurança.

\section{Requisitos Funcionais}
\begin{itemize}

    \item \textbf{RF01:} O sistema deve permitir a autenticação
    dos usuários.

    \item \textbf{RF02:} O sistema deve permitir o cadastro de
    motoristas secundários.

    \item \textbf{RF03:} O sistema deve permitir o envio de uma
    fotografia da CNH durante o cadastro.

    \item \textbf{RF04:} O sistema deve realizar a validação do
    CPF informado pelo usuário.

    \item \textbf{RF05:} O sistema deve diferenciar veículos
    novos e usados.

    \item \textbf{RF06:} O sistema deve permitir que o cliente
    consulte os veículos disponíveis para aluguel.

    \item \textbf{RF07:} O sistema deve permitir que o cliente
    realize reservas de veículos.

    \item \textbf{RF08:} O sistema deve calcular o valor do
    aluguel de acordo com as condições da reserva.

    \item \textbf{RF09:} O sistema deve permitir pagamentos por
    PIX e Cartão de Crédito.

    \item \textbf{RF10:} O sistema deve permitir a retirada e a
    devolução do veículo em locais diferentes da agência.

    \item \textbf{RF11:} O sistema deve calcular um adicional
    relacionado à devolução fora da agência.

    \item \textbf{RF12:} O sistema deve permitir o registro de
    fotografias do veículo durante a retirada e a devolução.

    \item \textbf{RF13:} O sistema deve permitir o registro da
    quilometragem do veículo.

    \item \textbf{RF14:} O sistema deve permitir o registro do
    nível de combustível do veículo.

    \item \textbf{RF15:} O sistema deve detectar automaticamente
    a quilometragem apresentada em uma fotografia do painel.

    \item \textbf{RF16:} O sistema deve permitir o registro de
    problemas ou danos identificados no veículo.

    \item \textbf{RF17:} O sistema deve permitir que o usuário
    consulte seu histórico de aluguéis.

    \item \textbf{RF18:} O sistema deve manter o histórico de
    utilização dos veículos.

    \item \textbf{RF19:} O sistema deve disponibilizar um painel
    de manutenção para a agência.

    \item \textbf{RF20:} O sistema deve registrar os custos
    operacionais dos veículos.

    \item \textbf{RF21:} O sistema deve calcular uma estimativa
    de poluição associada à utilização do veículo.

    \item \textbf{RF22:} O sistema deve permitir o acionamento
    de emergência.

    \item \textbf{RF23:} O sistema deve registrar a localização
    do usuário durante o acionamento de uma emergência.

    \item \textbf{RF24:} O sistema deve permitir a seleção ou
    descrição do tipo de emergência.

    \item \textbf{RF25:} O sistema deve permitir o acionamento
    do seguro a partir de um botão de emergência.

    \item \textbf{RF26:} O sistema deve gerar um alerta quando a
    desvalorização do veículo atingir ou ultrapassar 15\% do
    valor FIPE de um ano para o outro.

    \item \textbf{RF27:} O sistema deve permitir o bloqueio
    temporário de usuários de acordo com as regras de utilização
    definidas pela agência.



\end{itemize}

% =========================================================
% 6. REQUISITOS NÃO FUNCIONAIS
% =========================================================

\section{Requisitos Não Funcionais}

\begin{itemize}

    \item \textbf{RNF01:} O sistema deve possuir uma interface
    responsiva, permitindo seu acesso por diferentes tamanhos
    de tela.

    \item \textbf{RNF02:} O sistema deve apresentar uma
    interface intuitiva e de fácil utilização.

    \item \textbf{RNF03:} O sistema deve utilizar autenticação
    para controlar o acesso às funcionalidades restritas.

    \item \textbf{RNF04:} As senhas dos usuários devem ser
    armazenadas de forma segura, utilizando mecanismos
    apropriados de criptografia ou hash.

    \item \textbf{RNF05:} O sistema deve implementar níveis de
    permissão de acordo com o perfil do usuário.

\end{itemize}

% Casos de Uso ----------------------------------------------------------------------

\chapter{Diagramas}
\label{sec:Diagrama}
\section{Diagrama de Caso de Uso}

Seguindo o diagrama abaixo, temos os principais casos de uso do sistema

\begin{figure}[!h]
	\centering
	\includegraphics[scale=1]{imagens/USO.png}
	\caption{Diagrama de Casos de Uso}
	\label{Diagrama de Casos de Uso}
\end{figure}
\newpage

\section{Diagrama UML}
No diagrama UML, cada uma das classes é projetada com seus respectivos atributos e métodos.

\begin{figure}[h!]
    \centering
    \includegraphics[width =\linewidth]{imagens/uml.png}
    \caption{Diagrama UML}
    \label{fig:diagrama_uml}
\end{figure}
\newpage

% =========================================================
% ARQUITETURA DO SISTEMA
% =========================================================

\chapter{Arquitetura do Sistema}

\section{Visão Geral}

O sistema será desenvolvido seguindo uma arquitetura baseada na
separação entre as camadas de apresentação, aplicação e
persistência de dados. Essa organização tem como objetivo
facilitar a manutenção, a evolução e a integração entre os
diferentes componentes do sistema.

A camada de apresentação será responsável pela interface com os
usuários e será desenvolvida utilizando React. Por meio dela,
clientes e funcionários da agência poderão acessar as
funcionalidades disponibilizadas pelo sistema, como consulta de
veículos, realização de reservas, acompanhamento de aluguéis,
pagamentos e gerenciamento da frota.

A camada de aplicação será desenvolvida em Python e será
responsável pelo processamento das requisições, autenticação dos usuários e gerenciamento
das operações relacionadas aos veículos, reservas, pagamentos,
manutenções e demais funcionalidades do sistema.

A persistência dos dados será realizada utilizando MongoDB,
armazenando informações relacionadas aos
usuários, veículos, reservas, aluguéis, pagamentos, inspeções,
manutenções e demais registros do sistema.

A comunicação entre o front-end e o back-end será realizada por
meio de uma API, permitindo que a interface desenvolvida em
React envie requisições ao servidor e receba os dados necessários
para apresentar as informações aos usuários.

A arquitetura geral do sistema pode ser representada da
seguinte forma:

\begin{center}

\begin{tikzpicture}[
    node distance=3cm,
    every node/.style={
        draw,
        rounded corners,
        align=center,
        minimum width=4cm,
        minimum height=1.2cm
    }
]

\node (frontend) {Front-end\\React};

\node (backend) [right=of frontend] {Back-end\\Python};

\node (database) [right=of backend] {Banco de Dados\\MongoDB};

\draw[<->, thick] (frontend) -- 
node[above,draw=none] {API} (backend);

\draw[<->, thick] (backend) -- 
node[above,draw=none] {Dados} (database);

\end{tikzpicture}

\end{center}

\section{Fluxo de Comunicação}

O fluxo de comunicação do sistema ocorrerá principalmente entre
o front-end, o back-end e o banco de dados. Quando o usuário
realizar uma ação na interface, o front-end enviará uma
requisição para a API disponibilizada pelo back-end.

O back-end será responsável por receber a requisição, verificar
as informações fornecidas, aplicar as regras de negócio
correspondentes e, quando necessário, realizar operações no
banco de dados.

Após o processamento da requisição, o back-end retornará uma
resposta para o front-end, normalmente utilizando o formato
JSON. O front-end utilizará essas informações para atualizar a
interface e apresentar o resultado ao usuário.

De forma simplificada, o fluxo pode ser representado pelas
seguintes etapas:

\begin{enumerate}

    \item O usuário realiza uma ação na interface do sistema;

    \item O front-end React identifica a ação e envia uma
    requisição para a API;

    \item O back-end Python recebe e processa a requisição;

    \item Quando necessário, o back-end realiza operações de
    consulta, inserção, alteração ou exclusão no banco;

    \item O banco retorna os dados solicitados ou confirma a
    operação realizada;

    \item O back-end processa o resultado e envia uma resposta
    para o front-end;

    \item O React recebe a resposta e atualiza a interface
    apresentada ao usuário.

\end{enumerate}

Por exemplo, durante uma reserva de veículo, o fluxo poderá
ocorrer da seguinte maneira:

\begin{center}

\begin{tikzpicture}[
    node distance=1.5cm,
    every node/.style={
        draw,
        rounded corners,
        align=center,
        minimum width=4.2cm,
        minimum height=1cm
    }
]

\node (a) {Usuário\\realiza reserva};

\node (b) [below=of a] {Front-end\\React};

\node (c) [below=of b] {API\\Back-end Python};

\node (d) [below=of c] {MongoDB\\verificação de disponibilidade};

\node (e) [below=of d] {Back-end\\processamento da reserva};

\node (f) [below=of e] {Front-end\\confirmação da reserva};

\draw[->, thick] (a) -- (b);
\draw[->, thick] (b) -- (c);
\draw[->, thick] (c) -- (d);
\draw[->, thick] (d) -- (e);
\draw[->, thick] (e) -- (f);

\end{tikzpicture}

\end{center}

Dessa forma, a separação entre as camadas permite que cada
componente possua responsabilidades específicas, reduzindo o
acoplamento entre a interface, as regras de negócio e o
armazenamento dos dados.
% Tecnologias Utilizadas -------------------------------------------------------------
\chapter{Tecnologias Utilizadas}

\begin{itemize}
	\item \textbf{Back-end:} A linguagem escolhida para o desenvolvimento foi a linguagem Python.
	\item \textbf{Front-end:} O front-end servirá como ferramenta de visualização para o sistema e será desenvolvido principalmente com React.
	\item \textbf{Banco de Dados:} MongoDB
	\item \textbf{Controle de Versão:} Todo o processo será documentado e versionado dentro da plataforma github, que poderá ser acessada através do \href{https://github.com/DaviRodrish/POODEV2026}{Repositório Github}
\end{itemize}

% Calendário -------------------------------------------------------------------------
\chapter{Cronograma}
\label{sec:cronograma}


\vspace{2cm}

\begin{figure}[!h]
\centering
\includegraphics[width = \linewidth]{imagens/Calendario.png}
\caption{Cronograma de Desenvolvimento}
\label{Cronograma do Primeiro Bimestre de Desenvolvimento}
\end{figure}


\chapter{Orçamento}

No orçamento abaixo foram levados em consideração o período de três meses de desenvolvimento do sistema, somado com o periodo de uso de um ano seguido.
\begin{figure}[!h]
	\centering
	\includegraphics[width = \linewidth]{imagens/orcamento.png}
	\label{Orçamento}
	\caption{Orçamento}
\end{figure}

\end{document}

